\clearpage
\renewcommand{\sectioncolor}{sand}
\section{The 38 contigs: what coverage already tells us about copy number}
\markboth{Repeats and copy number}{}

\noindent\textbf{\color{sand}The biology question:} \emph{why does a single circular chromosome come
back as 38 pieces, and how many copies of each repeat are really there?}

\subsection{The answer is sitting in the contig names}

SPAdes writes the $k$-mer coverage into every contig header. Our eight large contigs
($>100$ kb) average \textbf{29.35}. Divide any contig's coverage by that number and you get its
\textbf{copy number}.

\begin{center}
\begin{tikzpicture}
\begin{axis}[width=15.4cm,height=6.4cm,
  xlabel={\footnotesize contig length (bp, log scale)},
  ylabel={\footnotesize coverage $\div$ 29.35 = copies},
  xmode=log, xmin=140,xmax=700000, ymin=0.5, ymax=9.2,
  ytick={1,2,3,4,5,6,7,8}, tick label style={font=\scriptsize}, label style={font=\footnotesize},
  grid=major, grid style={slate!18,line width=0.3pt}, axis line style={slate!60},
  clip=false]
 \addplot[only marks,mark=*,mark size=2.0pt,color=bio] coordinates
  {(479695,0.869)(392086,0.999)(357380,1.131)(319865,1.088)(296338,1.211)(293857,0.852)
   (235429,0.894)(169873,0.953)(62040,1.026)(39185,1.060)(37993,1.219)(35067,1.041)
   (20291,1.166)(15145,1.114)(13219,1.001)(1937,2.221)(683,1.372)(634,1.302)
   (411,1.472)(332,1.070)(303,1.260)(278,1.169)(243,0.286)(237,0.255)(171,1.664)};
 \addplot[only marks,mark=square*,mark size=2.6pt,color=alert] coordinates
  {(19885,3.307)(3121,6.845)(2600,8.268)(1610,6.107)(625,2.077)(553,2.038)(407,2.117)
   (380,2.765)(351,3.781)(351,1.483)(324,3.390)(233,6.829)(181,3.018)};
 \draw[slate,dashed,line width=0.7pt] (axis cs:140,1)--(axis cs:700000,1);
 \node[font=\scriptsize,text=slate,anchor=south west] at (axis cs:150,1.04) {single copy};
 \node[font=\scriptsize\bfseries,text=alert,anchor=west] at (axis cs:3400,8.30) {NODE\_18 $\approx$ 8 copies (2.6 kb)};
 \node[font=\scriptsize\bfseries,text=alert,anchor=west] at (axis cs:4200,6.85) {NODE\_17 $\approx$ 7 copies (3.1 kb)};
 \node[font=\scriptsize\bfseries,text=alert,anchor=west] at (axis cs:24000,3.31) {NODE\_14 $\approx$ 3 copies (19.9 kb)};
 \node[font=\scriptsize,text=bio,anchor=east] at (axis cs:420000,0.80) {the 8 big contigs};
\end{axis}
\end{tikzpicture}
\end{center}
\vspace{2pt}
\begin{center}\begin{minipage}{0.93\textwidth}\footnotesize\color{slate}
\textbf{Figure 6.}\; Every contig in \texttt{spades\_assembly/contigs.fasta}, plotted by length and
inferred copy number. Green circles: single copy. Red squares: collapsed repeats. The two contigs
sitting below 0.3 (NODE\_34, NODE\_35) are low-coverage fragments, probably contaminant or chimeric.
\end{minipage}\end{center}

\begin{labbox}[title={What is actually missing from our assembly}]
\begin{center}\small
\begin{tabular}{@{}l r r r@{}}
\toprule
\rowcolor{sand!20}
& \textbf{stored in the FASTA} & \textbf{present in the cell} & \textbf{difference}\\
\midrule
repeat sequence & 32{,}207 bp & $\approx$ 124{,}696 bp & \textbf{$\approx$ 92 kb collapsed}\\
\bottomrule
\end{tabular}
\end{center}
\vspace{3pt}
Our assembly totals 2{,}803{,}313 bp. A closed chromosome of this lineage is typically $\approx$ 2.75--2.85 Mb
\emph{including} all repeat copies. The $\approx$92 kb we are storing once instead of 6--8 times is
exactly why the assembly is in 38 pieces: \textbf{the assembler cannot tell which copy connects to
which neighbour.}
\end{labbox}

\subsection{Counting copies does not need 71$\times$}

\begin{center}\small
\begin{tabular}{@{}l r r r r r@{}}
\toprule
\rowcolor{sand!20}
\textbf{Contig} & \textbf{Length} & \textbf{Copies} & \textbf{CV at 71.6$\times$} & \textbf{CV at 20$\times$} & \textbf{CV at 10$\times$}\\
\midrule
NODE\_14 & 19{,}885 & 3 & 0.74 \% & 1.40 \% & 1.98 \%\\
NODE\_17 & 3{,}121 & 7 & 1.22 \% & 2.31 \% & 3.27 \%\\
NODE\_18 & 2{,}600 & 8 & 1.25 \% & 2.37 \% & 3.35 \%\\
NODE\_20 & 1{,}610 & 6 & 1.84 \% & 3.47 \% & 4.91 \%\\
\bottomrule
\end{tabular}
\end{center}
\vspace{-2pt}
\begin{center}\begin{minipage}{0.93\textwidth}\footnotesize\color{slate}
\textbf{Table 2.}\; Poisson counting precision, $\mathrm{CV}=1/\sqrt{k\,d\,L/\ell}$ with read length
$\ell=233$. To tell $k$ copies from $k+1$ you need CV well below $1/2k$; every row clears that at
10$\times$.
\end{minipage}\end{center}

\begin{dobox}[title={Action: report copy number with an interval, not a point}]
This directly addresses \textbf{Limitation 3} in our own \texttt{METHODS\_v1.7.md}: \emph{``the
$\varphi$Sa3 boundaries are inferred from gene positions, not from attL/attR detection. The
$\sim$42.4 kb span is an estimate.''}\\[5pt]
The coverage data supports a \emph{proper interval}. The right formulation is an
\textbf{integer-constrained estimate}: copy numbers must be whole numbers, and they must balance at
every junction of the assembly graph (what goes in must come out). Adding those two constraints to
the Poisson likelihood turns a guess into a confidence interval --- and it runs in seconds on a
graph this size.
\end{dobox}
